\documentclass[12pt, a4paper]{report}

% ─── Packages ────────────────────────────────────────────────────────────────
\usepackage[top=2.5cm, bottom=2.5cm, left=3cm, right=2.5cm]{geometry}
\usepackage{fontenc}
\usepackage{inputenc}
\usepackage{setspace}
\usepackage{titlesec}
\usepackage{graphicx}
\usepackage{caption}
\usepackage{subcaption}
\usepackage{float}
\usepackage{amsmath}
\usepackage{amssymb}
\usepackage{booktabs}
\usepackage{array}
\usepackage{multirow}
\usepackage{listings}
\usepackage{xcolor}
\usepackage{fancyhdr}
\usepackage{tocloft}
\usepackage{hyperref}
\usepackage{appendix}
\usepackage{cite}
\usepackage{tikz}
\usepackage{pgfplots}
\pgfplotsset{compat=1.18}

% ─── Hyperref setup ──────────────────────────────────────────────────────────
\hypersetup{
    colorlinks=true,
    linkcolor=black,
    citecolor=black,
    urlcolor=blue,
    pdftitle={DC Motor Cascade PID Control System},
    pdfauthor={Your Name},
}

% ─── Page style ──────────────────────────────────────────────────────────────
\pagestyle{fancy}
\fancyhf{}
\fancyhead[L]{\small\leftmark}
\fancyhead[R]{\small\thepage}
\fancyfoot{}
\renewcommand{\headrulewidth}{0.4pt}

% ─── Line spacing ────────────────────────────────────────────────────────────
\onehalfspacing

% ─── Chapter / section title format ─────────────────────────────────────────
\titleformat{\chapter}[display]
  {\normalfont\Large\bfseries}{\chaptertitlename\ \thechapter}{12pt}{\Large}
\titlespacing*{\chapter}{0pt}{-10pt}{20pt}

% ─── Code listing style ──────────────────────────────────────────────────────
\lstdefinestyle{cstyle}{
    language=C,
    basicstyle=\ttfamily\small,
    keywordstyle=\color{blue}\bfseries,
    commentstyle=\color{gray}\itshape,
    stringstyle=\color{orange},
    numbers=left,
    numberstyle=\tiny\color{gray},
    numbersep=8pt,
    frame=single,
    breaklines=true,
    tabsize=4,
    showstringspaces=false,
    captionpos=b,
}

\lstdefinestyle{vhdlstyle}{
    language=VHDL,
    basicstyle=\ttfamily\small,
    keywordstyle=\color{blue}\bfseries,
    commentstyle=\color{gray}\itshape,
    numbers=left,
    numberstyle=\tiny\color{gray},
    numbersep=8pt,
    frame=single,
    breaklines=true,
    tabsize=4,
    captionpos=b,
}

% ─── TOC depth ───────────────────────────────────────────────────────────────
\setcounter{tocdepth}{3}
\setcounter{secnumdepth}{3}

% ─────────────────────────────────────────────────────────────────────────────
\begin{document}

% ════════════════════════════════════════════════════════════════════════════
%  TITLE PAGE
% ════════════════════════════════════════════════════════════════════════════
\begin{titlepage}
    \centering
    \vspace*{2cm}

    {\Large \textbf{[University Name]}}\\[0.5cm]
    {\large [Department Name]}\\[2cm]

    \rule{\linewidth}{0.5pt}\\[0.5cm]
    {\LARGE \textbf{DC Motor Cascade PID Control System}}\\[0.3cm]
    {\large Based on STM32F407 and AX309 FPGA}\\[0.5cm]
    \rule{\linewidth}{0.5pt}\\[2cm]

    {\large
    \begin{tabular}{ll}
        \textbf{Author:}   & [Your Name] \\[0.3cm]
        \textbf{Student ID:} & [Your Student ID] \\[0.3cm]
        \textbf{Supervisor:} & [Supervisor Name] \\[0.3cm]
        \textbf{Module:}   & [Module Code / Name] \\[0.3cm]
        \textbf{Date:}     & \today \\
    \end{tabular}
    }

    \vfill
    {\small Academic Year 2025--2026}
\end{titlepage}

% ════════════════════════════════════════════════════════════════════════════
%  ABSTRACT
% ════════════════════════════════════════════════════════════════════════════
\chapter*{Abstract}
\addcontentsline{toc}{chapter}{Abstract}
\markboth{Abstract}{Abstract}

This report presents the design and implementation of a DC motor cascade
closed-loop control system using an STM32F407 microcontroller and an AX309
FPGA development board. The system achieves smooth speed regulation from
0 to 100 RPM via a cascade PID architecture comprising a speed outer loop
and a current inner loop.

Register-level C drivers were developed for all peripherals including UART,
ADC, SPI, I2C, and TIM3 quadrature encoder interface. The FPGA, programmed
in VHDL, implements a 12-bit PWM generator operating at approximately
12.2\,kHz, receiving duty cycle commands from the STM32 via SPI. Motor
current is measured by an INA240A2 bidirectional current shunt monitor and
fed back to the inner loop, while encoder-based speed feedback closes the
outer loop.

Experimental results confirm stable closed-loop operation with a steady-state
speed error below 1\%, PWM frequency accuracy of 12.207\,kHz (matching the
theoretical 50\,MHz/4096), and a no-load current measurement of approximately
30\,mA consistent with the motor datasheet.

\vspace{1em}
\noindent\textbf{Keywords:} cascade PID, STM32F407, FPGA, DC motor control,
current sensing, register-level programming, VHDL

% ════════════════════════════════════════════════════════════════════════════
%  ACKNOWLEDGMENTS
% ════════════════════════════════════════════════════════════════════════════
\chapter*{Acknowledgments}
\addcontentsline{toc}{chapter}{Acknowledgments}
\markboth{Acknowledgments}{Acknowledgments}

% TODO: 写致谢
[Write your acknowledgments here.]

% ════════════════════════════════════════════════════════════════════════════
%  TABLE OF CONTENTS / TABLES / FIGURES
% ════════════════════════════════════════════════════════════════════════════
\tableofcontents
\listoftables
\listoffigures

% ════════════════════════════════════════════════════════════════════════════
%  CHAPTER 1 — System Analysis and Design
% ════════════════════════════════════════════════════════════════════════════
\chapter{System Analysis and Design}

\section{Functional Requirements and Feasibility Analysis}

% TODO: 描述课程要求、系统需要实现的功能
The primary objective of this project is to implement a closed-loop DC motor
speed control system meeting the following requirements:

\begin{itemize}
    \item Speed range: 0--100\,RPM, user-adjustable via rotary potentiometer
    \item Steady-state speed error: $<1\%$
    \item Real-time display of speed, duty cycle, and motor current on OLED
    \item Serial debug output via UART at 115200\,baud
    \item Non-volatile parameter storage using SPI Flash (W25Q64)
\end{itemize}

\section{System Architecture Design}

The system adopts a layered architecture as shown in Figure~\ref{fig:arch}.

\begin{figure}[H]
    \centering
    % TODO: 插入系统框图
    % \includegraphics[width=0.85\textwidth]{figures/system_arch.png}
    \fbox{\parbox{0.85\textwidth}{\centering\vspace{2cm}
    [System block diagram: Sensing → Decision → Actuation]
    \vspace{2cm}}}
    \caption{System architecture of the DC motor cascade PID control system.}
    \label{fig:arch}
\end{figure}

The data flow is:
\[
\text{Potentiometer} \xrightarrow{\text{ADC}} \text{STM32 PID}
\xrightarrow{\text{SPI2}} \text{FPGA} \xrightarrow{\text{PWM}}
\text{TB6612} \rightarrow \text{Motor}
\]
with encoder speed and INA240 current feeding back directly to the STM32.

\section{Control Strategy: Why Cascade PID}

A single-loop proportional--integral--derivative (PID) controller acting
directly on speed exhibits slow disturbance rejection when load current
changes. By adding an inner current loop, the current is regulated at a much
higher bandwidth, effectively linearising the plant for the outer speed loop
and improving overall dynamic performance.

The two loops operate at the same 1\,kHz rate but with separate tuning:
the inner loop uses a higher proportional gain to ensure it settles within
one outer-loop period.

% ════════════════════════════════════════════════════════════════════════════
%  CHAPTER 2 — Hardware Development
% ════════════════════════════════════════════════════════════════════════════
\chapter{Hardware Development}

\section{Microcontroller System (STM32F407G-DISC1)}

\subsection{MCU Selection and Specifications}

The STM32F407VG (Cortex-M4, 168\,MHz, FPU) was selected for its rich
peripheral set and sufficient processing power to run a 1\,kHz control
interrupt with $<1\%$ CPU load. Key peripherals used are listed in
Table~\ref{tab:mcu_periph}.

\begin{table}[H]
    \centering
    \caption{STM32F407 peripherals used in this project.}
    \label{tab:mcu_periph}
    \begin{tabular}{lll}
        \toprule
        \textbf{Peripheral} & \textbf{Function} & \textbf{Clock} \\
        \midrule
        TIM6     & 1\,kHz PID control interrupt & APB1 84\,MHz \\
        TIM3     & Quadrature encoder interface & APB1 84\,MHz \\
        ADC1     & Potentiometer + current sensing & APB2 21\,MHz \\
        SPI1     & W25Q64 Flash (not yet wired) & APB2 5.25\,MHz \\
        SPI2     & FPGA duty cycle transmission & APB1 5.25\,MHz \\
        I2C1     & OLED SSD1306 display & APB1 100\,kHz \\
        USART2   & Serial debug (CH340G) & APB1 115200 baud \\
        \bottomrule
    \end{tabular}
\end{table}

\subsection{Pin Allocation and Peripheral Mapping}

Table~\ref{tab:pins} lists all 16 signal pins used.

\begin{table}[H]
    \centering
    \caption{STM32F407 pin assignment.}
    \label{tab:pins}
    \begin{tabular}{lll}
        \toprule
        \textbf{Pin} & \textbf{Peripheral} & \textbf{Connected to} \\
        \midrule
        PA0  & ADC1\_IN0   & Potentiometer (via RC filter) \\
        PA1  & ADC1\_IN1   & INA240A2 output \\
        PA2  & USART2\_TX  & CH340G RXD \\
        PA3  & USART2\_RX  & CH340G TXD \\
        PA4  & SPI1\_CS    & W25Q64 chip select \\
        PA5  & SPI1\_SCK   & W25Q64 clock \\
        PA6  & SPI1\_MISO  & W25Q64 data out \\
        PA7  & SPI1\_MOSI  & W25Q64 data in \\
        PB6  & I2C1\_SCL   & OLED SCL \\
        PB7  & I2C1\_SDA   & OLED SDA \\
        PB12 & SPI2\_CS    & FPGA cs\_n \\
        PB13 & SPI2\_SCK   & FPGA sck \\
        PB14 & SPI2\_MISO  & FPGA miso (not used) \\
        PB15 & SPI2\_MOSI  & FPGA mosi \\
        PC6  & TIM3\_CH1   & Encoder channel A \\
        PC7  & TIM3\_CH2   & Encoder channel B \\
        \bottomrule
    \end{tabular}
\end{table}

\section{FPGA PWM Generation Module (AX309 Spartan-6)}

\subsection{Role of the FPGA}

The FPGA operates as a dedicated PWM generator. It receives a 12-bit duty
cycle value from the STM32 via SPI and outputs a corresponding PWM signal
to the TB6612FNG motor driver. This offloads real-time PWM generation from
the MCU and allows a higher PWM frequency than the STM32 timer could
conveniently achieve.

\subsection{PWM Frequency Derivation}

The AX309 board provides a 50\,MHz on-board oscillator. A 12-bit up-counter
resets at 4096, giving:

\begin{equation}
    f_{\text{PWM}} = \frac{f_{\text{clk}}}{2^{12}} = \frac{50 \times 10^6}{4096}
    \approx 12.207\,\text{kHz}
    \label{eq:pwm_freq}
\end{equation}

This frequency is above the audible range and provides sufficient resolution
(one LSB $= 100\,\%/4096 \approx 0.024\,\%$) for smooth speed control.
Experimentally verified with a logic analyser: 12.207\,kHz.

\section{Motor Drive Circuit (TB6612FNG)}

% TODO: 描述 TB6612 H桥原理，AIN1/AIN2/STBY 接线
The TB6612FNG dual H-bridge driver converts the 3.3\,V FPGA PWM signal to
the current levels required by the JGA25-370 DC motor. In this project the
motor runs in a single direction (AIN1 = HIGH, AIN2 = LOW, STBY = HIGH,
all hardwired), and speed is controlled solely by the PWM duty cycle applied
to the PWMA input.

\begin{table}[H]
    \centering
    \caption{TB6612FNG control signal wiring.}
    \label{tab:tb6612}
    \begin{tabular}{lll}
        \toprule
        \textbf{TB6612 pin} & \textbf{Connected to} & \textbf{Note} \\
        \midrule
        PWMA & FPGA J3-31 (F16) & PWM duty from FPGA \\
        AIN1 & 3.3\,V           & Direction: forward \\
        AIN2 & GND              & Direction: forward \\
        STBY & 3.3\,V           & Must be HIGH to enable \\
        VM   & 12\,V supply     & Motor power \\
        AO1/AO2 & Motor terminals & Output \\
        \bottomrule
    \end{tabular}
\end{table}

\section{Sensing System}

\subsection{Speed Sensing: Encoder and TIM3 Quadrature Decoding}

The JGA25-370 motor has a built-in Hall-effect quadrature encoder. TIM3 is
configured in Encoder Mode 3 (both edges of both channels), giving a
$\times4$ resolution multiplier. The speed is computed every 1\,ms:

\begin{equation}
    \text{RPM} = \frac{\Delta_{\text{count}} \times 60{,}000}
                      {N_{\text{PPR}}}
    \label{eq:rpm}
\end{equation}

where $\Delta_{\text{count}}$ is the signed 16-bit counter difference (handles
overflow automatically) and $N_{\text{PPR}}$ is the encoder pulses per
revolution (including gear ratio and $\times4$ mode).

An 8-point moving average filter smooths the raw RPM, and a hardware digital
filter (IC1F = IC2F = 0001, 2-sample) rejects spikes shorter than two
sampling periods.

\subsection{Current Sensing: INA240A2 Signal Chain}

\begin{figure}[H]
    \centering
    % TODO: 插入 INA240 信号链框图
    \fbox{\parbox{0.7\textwidth}{\centering\vspace{1.5cm}
    [INA240A2 signal chain diagram]
    \vspace{1.5cm}}}
    \caption{INA240A2 current sensing signal chain.}
    \label{fig:ina240}
\end{figure}

The INA240A2 (gain = 50\,V/V) monitors the voltage across a 0.1\,$\Omega$
shunt resistor. Its output voltage is:

\begin{equation}
    V_{\text{out}} = V_{\text{ref}} + G \cdot I \cdot R_{\text{shunt}}
    \label{eq:ina240_out}
\end{equation}

where $V_{\text{ref}} = V_{\text{CC}}/2$. The empirically calibrated
zero-current reference is $V_{\text{ref}} = 1773\,\text{mV}$ (ADC raw
$\approx 2200$). The motor current is recovered as:

\begin{equation}
    I\,[\text{mA}] = \frac{V_{\text{out}}\,[\text{mV}] - 1773}
                         {50 \times 0.1}
                  = \frac{V_{\text{out}} - 1773}{5}
    \label{eq:current}
\end{equation}

Validation: at no load and 100\,RPM, the measured current is approximately
30\,mA, consistent with the JGA25-370 datasheet specification.

\subsection{Speed Setpoint: Potentiometer and RC Low-Pass Filter}

A rotary potentiometer connected to PA0 (ADC1\_IN0) provides the target
speed reference. An RC low-pass filter ($f_c = 1/(2\pi RC)$) is inserted
between the wiper and the ADC pin to attenuate power-supply noise. Software
applies a 16-point sliding average and a dead zone (raw ADC $< 50
\Rightarrow$ setpoint $= 0$) to eliminate drift near zero.

% ════════════════════════════════════════════════════════════════════════════
%  CHAPTER 3 — Embedded Software Development
% ════════════════════════════════════════════════════════════════════════════
\chapter{Embedded Software Development}

All drivers are implemented at the register level (no HAL peripheral
abstractions) to demonstrate direct hardware control and satisfy course
requirements.

\section{Register-level Peripheral Drivers}

\subsection{UART — Serial Debug (USART2)}

% TODO: 说明 BRR 寄存器计算，AF7 配置，CH340 接线
USART2 is configured for 115200\,baud, 8N1. With APB1 = 42\,MHz:
\begin{equation}
    \text{BRR} = \frac{42{,}000{,}000}{16 \times 115{,}200} \approx 22.786
    \quad\Rightarrow\quad \texttt{0x16D}
\end{equation}
The ST-Link VCP was found non-functional on this board; an external CH340G
USB-to-TTL module (PA2/PA3) is used instead.

\subsection{ADC — Dual-Channel Acquisition}

% TODO: ADCPRE、采样时间、软件触发流程
ADC1 uses software triggering with dynamic channel selection (SQR3). The
clock is divided from APB2 (84\,MHz) by 4 to give 21\,MHz. Sampling times
are set to 84 cycles for the high-impedance potentiometer channel (IN0) and
28 cycles for the low-impedance INA240 output (IN1).

\subsection{SPI1 — W25Q64 Flash Storage}

% TODO: W25Q64 指令集，ReadID，Page Program，Sector Erase
[Describe SPI1 configuration and W25Q64 command protocol. Note: chip not
yet physically wired; driver validated via SPI1 loopback test (PASS).]

\subsection{SPI2 — FPGA Communication Protocol}

Data is sent to the FPGA as two bytes encoding the 12-bit duty cycle:

\begin{verbatim}
CS↓
  Byte 0: [7:4]=0000  [3:0]=duty[11:8]   (high nibble)
  Byte 1: [7:0]=duty[7:0]                (low byte)
CS↑
\end{verbatim}

The CS edge serves as the frame delimiter; no synchronisation header is
required. SPI2 is configured for APB1/8 = 5.25\,MHz, Mode 0, 8-bit MSB-first.
Since the FPGA does not return data, MISO (PB14) may be left unconnected.

\subsection{I2C1 — OLED SSD1306 Display}

% TODO: CCR/TRISE 计算，100kHz，OLED 初始化序列
I2C1 operates at 100\,kHz. The CCR value is set to 180 with TRISE = 37 for
APB1 = 36\,MHz. A timeout counter and flag-clearing sequence (AF, BERR,
ARLO, OVR) are applied after each transaction to prevent the bus from
locking up.

\subsection{TIM3 — Quadrature Encoder Interface}

% TODO: SMCR SMS=011，CCMR1 滤波器，CCER 极性
TIM3 is set to Encoder Mode 3 (SMS[2:0] = 011), where both rising and
falling edges of both channels increment or decrement the 16-bit counter.
A 2-sample digital filter (IC1F = IC2F = 0001) is enabled to reject noise.

\section{Signal Conditioning}

\subsection{Moving Average Filters}

\begin{table}[H]
    \centering
    \caption{Software filters applied to each signal.}
    \label{tab:filters}
    \begin{tabular}{llll}
        \toprule
        \textbf{Signal} & \textbf{Window} & \textbf{Delay} & \textbf{Function} \\
        \midrule
        Potentiometer ADC & 16 samples & 16\,ms & \texttt{ADC1\_Read\_Filtered()} \\
        Encoder RPM       & 8 samples  & 8\,ms  & \texttt{RPM\_Filter()} \\
        Motor current     & 8 samples  & 8\,ms  & \texttt{ADC1\_ReadCurrent\_Filtered\_mA()} \\
        \bottomrule
    \end{tabular}
\end{table}

\subsection{Dead Zone and Anti-windup}

A dead zone of ADC $< 50$ forces \texttt{setpoint\_rpm = 0} to suppress
integrator wind-up caused by the potentiometer's natural offset near
zero. Integrator clamping (\texttt{integral\_max}) is applied inside
\texttt{PID\_Calc()} to both loops independently.

% ════════════════════════════════════════════════════════════════════════════
%  CHAPTER 4 — Cascade PID Control Algorithm
% ════════════════════════════════════════════════════════════════════════════
\chapter{Cascade PID Control Algorithm}

\section{PID Mathematical Model}

The positional discrete PID controller is derived from the continuous form:

\begin{equation}
    u(t) = K_p\,e(t) + K_i\int_0^t e(\tau)\,d\tau + K_d\,\frac{de(t)}{dt}
\end{equation}

Discretised with sampling period $T$ (here $T = 1\,\text{ms}$):

\begin{equation}
    u[n] = K_p\,e[n]
         + K_i \sum_{k=0}^{n} e[k]
         + K_d\,\bigl(e[n] - e[n-1]\bigr)
    \label{eq:pid_discrete}
\end{equation}

Integrator wind-up is prevented by clamping the accumulated sum:
\begin{equation}
    \sum e[k] \;\leftarrow\; \text{clamp}\!\left(\sum e[k],\;
    -I_{\max},\; +I_{\max}\right)
\end{equation}

and the output is similarly bounded to $[u_{\min},\, u_{\max}]$.

\section{Cascade Structure Design}

\begin{figure}[H]
    \centering
    % TODO: 画串级 PID 框图
    \fbox{\parbox{0.85\textwidth}{\centering\vspace{2cm}
    [Cascade PID block diagram: speed loop → current loop → plant]
    \vspace{2cm}}}
    \caption{Cascade PID control structure.}
    \label{fig:cascade}
\end{figure}

\begin{itemize}
    \item \textbf{Speed outer loop}: $\omega_{\text{ref}} \to I_{\text{ref}}$
    (RPM setpoint from potentiometer $\to$ current demand in mA).
    \item \textbf{Current inner loop}: $I_{\text{ref}} \to d$
    (current demand $\to$ 12-bit PWM duty cycle 0--4095).
\end{itemize}

For stable cascade operation the inner loop bandwidth must be at least
3--5$\times$ higher than the outer loop~\cite{astrom2006}. Both loops run
at 1\,kHz here; a higher inner-loop gain achieves the required relative
bandwidth.

\section{Speed Outer Loop Design}

\begin{table}[H]
    \centering
    \caption{Speed outer loop PID parameters.}
    \label{tab:speed_pid}
    \begin{tabular}{lll}
        \toprule
        \textbf{Parameter} & \textbf{Value} & \textbf{Rationale} \\
        \midrule
        $K_p$ & 1.5  & 100\,RPM error $\Rightarrow$ P term = 150\,mA \\
        $K_i$ & 0.1  & Integrator converges in $\sim$1\,s \\
        $K_d$ & 0.0  & Encoder noise renders derivative impractical \\
        $u_{\max}$ & 400\,mA & Measured motor operating range: 200--300\,mA \\
        $I_{\max}$ & 4000    & $= u_{\max}/K_i$ \\
        \bottomrule
    \end{tabular}
\end{table}

\section{Current Inner Loop Design}

\begin{table}[H]
    \centering
    \caption{Current inner loop PID parameters.}
    \label{tab:current_pid}
    \begin{tabular}{lll}
        \toprule
        \textbf{Parameter} & \textbf{Value} & \textbf{Rationale} \\
        \midrule
        $K_p$ & 5.0  & Reduced from 20 after limit-cycle oscillation observed \\
        $K_i$ & 0.2  & Integrates slowly to avoid overshoot \\
        $K_d$ & 0.0  & Current signal is relatively clean \\
        $u_{\max}$ & 4095 & Full 12-bit PWM range \\
        $I_{\max}$ & 8000 & $\approx u_{\max}/K_i$ \\
        \bottomrule
    \end{tabular}
\end{table}

\section{1\,kHz Control Loop Implementation}

The \texttt{TIM6\_DAC\_IRQHandler} executes every 1\,ms in the following
sequence:

\begin{enumerate}
    \item Update encoder: compute $\Delta_{\text{count}}$ and filtered RPM.
    \item Read potentiometer ADC (16-point filtered); apply dead zone.
    \item \textbf{Speed PID}: $e_\omega = \omega_{\text{ref}} - \omega_{\text{meas}}
          \;\Rightarrow\; I_{\text{ref}}$.
    \item Read current ADC (8-point filtered).
    \item \textbf{Current PID}: $e_I = I_{\text{ref}} - I_{\text{meas}}
          \;\Rightarrow\; d \in [0, 4095]$.
    \item Transmit duty $d$ to FPGA via SPI2 (2 bytes, $\approx3\,\mu$s).
\end{enumerate}

Total ISR execution time $\approx 9\,\mu\text{s}$ (0.9\% of the 1\,ms period).

% ════════════════════════════════════════════════════════════════════════════
%  CHAPTER 5 — FPGA Development
% ════════════════════════════════════════════════════════════════════════════
\chapter{FPGA Development}

The FPGA design consists of three VHDL modules synthesised with Xilinx ISE
14.7 for the Spartan-6 XC6SLX9-2FTG256 device.

\section{SPI Slave Receiver (\texttt{spi\_slave.vhd})}

\subsection{Cross-clock-domain Synchronisation}

The SPI signals (SCK, CS\_N, MOSI) are asynchronous to the 50\,MHz FPGA
clock. A two-stage synchronisation chain eliminates metastability:

\begin{lstlisting}[style=vhdlstyle, caption={Two-stage synchroniser for SCK.}]
p_sync : process(clk)
begin
    if rising_edge(clk) then
        sck_d1 <= sck;
        sck_d2 <= sck_d1;
    end if;
end process;
sck_rise <= sck_d1 and (not sck_d2);  -- single-cycle rise pulse
\end{lstlisting}

\subsection{Receive State Machine}

Data is shifted in on each detected SCK rising edge while CS\_N is low.
After 8 bits, the high nibble is latched; after a second 8 bits the full
12-bit duty cycle is assembled and registered on the CS\_N rising edge.

\section{PWM Generator (\texttt{pwm\_gen.vhd})}

A 12-bit free-running counter compares against the latched duty register:

\begin{lstlisting}[style=vhdlstyle, caption={PWM comparator.}]
pwm_out <= '1' when (pwm_cnt < duty_reg) else '0';
\end{lstlisting}

The \texttt{duty\_valid} strobe from the SPI slave updates \texttt{duty\_reg}
only at frame boundaries, preventing glitches mid-period.

\section{Top-level Integration (\texttt{top.vhd})}

\texttt{top.vhd} instantiates \texttt{spi\_slave} and \texttt{pwm\_gen} and
routes the duty bus between them. The pin constraints are specified in
\texttt{pin.ucf}; Table~\ref{tab:fpga_pins} summarises the FPGA pin
assignments.

\begin{table}[H]
    \centering
    \caption{FPGA pin assignment (AX309 J3 expansion connector).}
    \label{tab:fpga_pins}
    \begin{tabular}{llll}
        \toprule
        \textbf{Signal} & \textbf{FPGA pin} & \textbf{J3 pin} & \textbf{Connected to} \\
        \midrule
        clk     & P8  & On-board oscillator & --- \\
        sck     & C15 & J3-21 & STM32 PB13 \\
        cs\_n   & D16 & J3-23 & STM32 PB12 \\
        mosi    & C16 & J3-24 & STM32 PB15 \\
        miso    & E15 & J3-26 & STM32 PB14 (not used) \\
        pwm\_out & F16 & J3-31 & TB6612 PWMA \\
        \bottomrule
    \end{tabular}
\end{table}

% ════════════════════════════════════════════════════════════════════════════
%  CHAPTER 6 — System Integration and Testing
% ════════════════════════════════════════════════════════════════════════════
\chapter{System Integration and Testing}

\section{Debugging Process and Lessons Learned}

Table~\ref{tab:issues} summarises the significant issues encountered during
development and their resolutions.

\begin{table}[H]
    \centering
    \caption{Debugging history: problems encountered and resolved.}
    \label{tab:issues}
    \begin{tabular}{p{3cm}p{4cm}p{5.5cm}}
        \toprule
        \textbf{Problem} & \textbf{Symptom} & \textbf{Root cause \& fix} \\
        \midrule
        ST-Link VCP no output &
            UART silent on COM6 &
            ST-Link internal UART sub-module inoperative; switched to
            external CH340G on PA2/PA3. \\
        \addlinespace
        INA240 gain mismatch &
            Current reading $2.5\times$ too high &
            Module is A2 (50\,V/V), not A1 (20\,V/V); corrected gain
            constant. \\
        \addlinespace
        Current-loop limit-cycle oscillation &
            Duty oscillates 0--100\% &
            $K_p=20$ too aggressive; reduced to 5, $K_i$ from 0.5 to 0.2. \\
        \addlinespace
        OLED lockup &
            Screen freezes after first frame &
            I2C PE disabled mid-transaction; added timeout guard and
            cleared AF/BERR/ARLO/OVR flags. \\
        \addlinespace
        RPM = 0, duty saturated &
            Motor appears stationary &
            Encoder 5\,V supply disconnected; speed loop integrator
            winds up; fix: verify encoder power first. \\
        \addlinespace
        ADC drift at zero &
            Motor twitches at rest &
            Potentiometer offset $\approx9$ ADC counts; dead zone
            (ADC $<50 \Rightarrow$ setpoint $=0$) applied. \\
        \bottomrule
    \end{tabular}
\end{table}

\section{Experimental Results}

\subsection{PWM Frequency Verification}

The FPGA PWM output on J3-31 was measured with a logic analyser.
The result of 12.207\,kHz matches the theoretical value of
$50\,\text{MHz}/4096 = 12.207\,\text{kHz}$ to within measurement
uncertainty.

\subsection{Speed Control Performance}

% TODO: 插入串口数据或实测曲线
Figure~\ref{fig:speed_response} shows the closed-loop speed response when
the potentiometer is stepped from 0 to full scale.

\begin{figure}[H]
    \centering
    \fbox{\parbox{0.75\textwidth}{\centering\vspace{3cm}
    [Speed response plot: RPM vs time]
    \vspace{3cm}}}
    \caption{Closed-loop speed step response (0 to 100\,RPM).}
    \label{fig:speed_response}
\end{figure}

Steady-state measurements at various setpoints are listed in
Table~\ref{tab:speed_results}.

\begin{table}[H]
    \centering
    \caption{Steady-state speed control results.}
    \label{tab:speed_results}
    \begin{tabular}{rrrr}
        \toprule
        \textbf{Setpoint (RPM)} & \textbf{Measured (RPM)} &
        \textbf{Error (RPM)} & \textbf{Error (\%)} \\
        \midrule
        20  & % TODO
            & & \\
        50  & & & \\
        80  & & & \\
        100 & & & \\
        \bottomrule
    \end{tabular}
\end{table}

\subsection{Current Sensing Accuracy}

At no load and 100\,RPM, the measured motor current is approximately
30\,mA, consistent with the JGA25-370 datasheet no-load current
specification (20--150\,mA).

\section{Conclusion and Future Work}

This project successfully demonstrated register-level implementation of a
cascade PID DC motor control system on an STM32F407 + FPGA platform.
Key achievements include:

\begin{itemize}
    \item Stable closed-loop speed regulation over 0--100\,RPM
    \item PWM generation accuracy confirmed at 12.207\,kHz
    \item Current sensing validated at no-load ($\approx30$\,mA)
    \item All register-level drivers written and verified
\end{itemize}

Remaining work includes:
\begin{itemize}
    \item Physical wiring and verification of the W25Q64 Flash chip
    \item Load testing of the current inner loop
    \item Optimisation of PID parameters under varying load conditions
\end{itemize}

% ════════════════════════════════════════════════════════════════════════════
%  REFERENCES
% ════════════════════════════════════════════════════════════════════════════
\bibliographystyle{IEEEtran}
\begin{thebibliography}{9}

\bibitem{rm0090}
STMicroelectronics,
\textit{RM0090: STM32F405/407/415/417 Reference Manual},
Rev.\ 19, 2021. [Online]. Available: \url{https://www.st.com}

\bibitem{ds8626}
STMicroelectronics,
\textit{DS8626: STM32F407VG Datasheet},
Rev.\ 9, 2021.

\bibitem{ina240}
Texas Instruments,
\textit{INA240 Bidirectional, Low- and High-Side, Voltage Output,
Current-Sense Amplifier Datasheet}, SBOS566D, 2015.

\bibitem{w25q64}
Winbond Electronics,
\textit{W25Q64FV 64M-bit Serial Flash Memory Datasheet},
Rev.\ L, 2017.

\bibitem{tb6612}
Toshiba,
\textit{TB6612FNG Dual DC Motor Driver IC Datasheet},
2015.

\bibitem{ssd1306}
Solomon Systech,
\textit{SSD1306 Dot Matrix OLED/PLED Segment/Common Driver with
Controller Datasheet}, 2008.

\bibitem{astrom2006}
K.~J. \AA{}str\"{o}m and T.~H\"{a}gglund,
\textit{Advanced PID Control}.
ISA -- The Instrumentation, Systems, and Automation Society, 2006.

\end{thebibliography}

% ════════════════════════════════════════════════════════════════════════════
%  APPENDIX
% ════════════════════════════════════════════════════════════════════════════
\appendix

\chapter{Key Source Code}

\section{TIM6 1\,kHz PID Control Loop}

\begin{lstlisting}[style=cstyle, caption={TIM6\_DAC\_IRQHandler: 1\,kHz cascade PID loop.}]
void TIM6_DAC_IRQHandler(void)
{
    if (TIM6->SR & TIM_SR_UIF) {
        TIM6->SR &= ~TIM_SR_UIF;

        Encoder_Update();

        uint16_t adc_val      = ADC1_Read_Filtered();
        float    setpoint_rpm = (adc_val < 50U) ? 0.0f
                             : (float)adc_val * (MOTOR_MAX_RPM / 4095.0f);

        float current_setpoint = PID_Calc(&speed_pid,
                                          setpoint_rpm, g_encoder_rpm);

        float measured_current = ADC1_ReadCurrent_Filtered_mA();

        float    duty_f = PID_Calc(&current_pid,
                                   current_setpoint, measured_current);
        uint16_t duty   = (uint16_t)duty_f;

        uint8_t hi = (duty >> 8U) & 0x0FU;
        uint8_t lo =  duty        & 0xFFU;
        FPGA_CS_LOW();
        SPI2_ReadWriteByte(hi);
        SPI2_ReadWriteByte(lo);
        FPGA_CS_HIGH();

        g_pid_duty = duty;
    }
}
\end{lstlisting}

\section{Positional PID Implementation}

\begin{lstlisting}[style=cstyle, caption={PID\_Calc: positional PID with anti-windup.}]
float PID_Calc(PID_TypeDef *pid, float setpoint, float feedback)
{
    float err = setpoint - feedback;

    pid->integral += err;
    if      (pid->integral >  pid->integral_max)
        pid->integral =  pid->integral_max;
    else if (pid->integral < -pid->integral_max)
        pid->integral = -pid->integral_max;

    float output = pid->kp * err
                 + pid->ki * pid->integral
                 + pid->kd * (err - pid->err_prev);
    pid->err_prev = err;

    if      (output > pid->out_max) output = pid->out_max;
    else if (output < pid->out_min) output = pid->out_min;
    return output;
}
\end{lstlisting}

\end{document}
